\section{Theoretische Grundlagen}

\subsection{Grundlagen des Supervised Learning}
Dieses Kapitel definiert das mathematische Fundament des Supervised Learning und erklärt Loss-Funktionen sowie Optimizer zur Fehlerminimierung während des Trainings.

\subsection{Behavior Cloning}

\subsubsection{Mathematische Formulierung}
Behavior Cloning wird als Imitation Learning in Form eines Regressionsproblems für ein direktes State-Action-Mapping formalisiert.

\subsubsection{Covariate Shift und Compounding Errors}
Es wird die Fehlerakkumulation erklärt, die auftritt, sobald das Modell von den im Training beobachteten States abweicht.

\subsection{Geeignete Modellarchitekturen}

\subsubsection{Multi-Layer Perceptrons}
Dieser Abschnitt erläutert Aufbau und Funktion einfacher Feedforward-Netze ohne zeitliches Gedächtnis.

\subsubsection{Recurrent Neural Networks}
Dieser Abschnitt beschreibt die Funktionsweise von grundlegenden RNNs zur Verarbeitung sequenzieller Zeitreihendaten mit rekurrenten Verbindungen.

\subsubsection{Long Short-Term Memory Networks}
LSTMs erweitern RNNs durch Speicherzellen und Gate-Mechanismen, um langfristige Abhängigkeiten zu erfassen und das Problem verschwindender Gradienten zu reduzieren.

\subsubsection{Gated Recurrent Units}
GRUs vereinfachen die LSTM-Architektur durch weniger Parameter bei ähnlicher Ausdruckskraft und ermöglichen schnelleres Training bei vergleichbarer Leistung.

\subsection{Overfitting-Prävention}

\subsubsection{Weight Initialization}
Es werden Methoden für einen stabilen Trainingsstart zur Vermeidung explodierender oder verschwindender Gradienten beschrieben.

\subsubsection{Regularisierungstechniken}
Dieser Abschnitt erklärt die mathematische Bestrafung von Gewichten (Weight Decay) und den Einsatz von Dropout-Layern.
